% ============================================================
% ch20a.tex —— 第 20 章「Seiberg–Witten 理论」前半（20.1–20.2.5）
% 印刷 pp.287–309（PDF p327–p349，23 页），图 20.1–20.6，编号公式 20.1–20.13
% 底本：Springer LNM 2386, Ch.20 Seiberg–Witten theory
% ============================================================

\chapter{Seiberg–Witten 理论}
% [原文 p.287]

\ChapterAbstract{Seiberg--Witten 理论给出 $\mathcal{N}=2$ 超对称规范理论在库仑分支上量子动力学的精确解，并揭示穿墙现象。这些进展使得能够自洽地实现 $\mathcal{N}=2$ 的 S 对偶，即使在 BPS 谱似乎缺乏必要对称性的情形中亦然。我们从四维 $\mathcal{N}=2$ $\mathfrak{su}(2)$ SYM 理论的 Seiberg--Witten 分析开始，既讨论纯理论，也讨论带至多四个基本超多重态的情形，然后概述它到其他 $\mathcal{N}=2$ 规范理论的推广。我们的讨论遵循原始文献 [313, 314, 114] 与重要综述 [18, 245, 106, 339]，更偏数学的视角则见 [367, 137]。}

\section{纯 $\mathfrak{su}(2)$ 超 Yang--Mills 理论}

纯 $4d$ $\mathcal{N}=2$ $\mathfrak{su}(2)$ SYM 的拉格朗日密度具有式 (18.3) 的形式。由命题 18.21，可以安全地把式 (18.3) 中的 $\mathcal{N}=2$ 矢量超场重标度为 $1/2g$，得到：
\begin{equation}
\begin{aligned}
\mathcal{L}&=\frac{1}{4\pi}\operatorname{Im}\Bigl(\tilde\tau\int\dd^2\theta\,\operatorname{tr}W^\alpha W_\alpha\Bigr)+\frac{1}{g^2}\operatorname{tr}\int\dd^2\theta\,\dd^2\bar\theta\,\Phi \ee^{[V,\cdot]}\Phi\\
&=\frac{1}{g^2}\operatorname{tr}\Bigl[-\frac{1}{2}F_{\mu\nu}F^{\mu\nu}-2i\lambda\sigma^\mu D_\mu\bar\lambda-2i\tilde\lambda\sigma^\mu D_\mu\bar{\tilde\lambda}+\frac{\Theta g^2}{16\pi^2}F_{\mu\nu}(\star F)^{\mu\nu}\\
&\qquad+2\bar D_\mu\bar\phi D^\mu\phi+D^2+2\bar FF+2i\sqrt2 g\,\bar\phi\{\lambda,\tilde\lambda\}-2i\sqrt2 g\{\bar\lambda,\bar{\tilde\lambda}\}\phi+2gD[\phi,\bar\phi]\Bigr],
\end{aligned}
\tag{20.1}
\end{equation}
其中复化规范耦合现在记作 $\tilde\tau$。它随标度对数地跑动：
\begin{equation}
\tilde\tau(\mu)=-\frac{b_1}{2\pi i}\log\Bigl(\frac{\mu}{\Lambda}\Bigr)-\frac{i}{\pi}\sum_k c_k\Bigl(\frac{\Lambda}{\mu}\Bigr)^{4k},
\tag{20.2}
\end{equation}
% [原文 p.288]
其中 $b_1=2N-F=4$（如例 18.19），$\Lambda$ 是式 (18.20) 的动力学标度，而非微扰修正的一般形式由全纯性所约束。在式 (20.2) 的求和中，每个 $c_k$ 是编码瞬子数为 $k$ 的规范场位形对 $\tilde\tau$ 跑动之贡献的复数。由于 $b_1>0$，纯 $4d$ $\mathcal{N}=2$ $\mathfrak{su}(2)$ SYM 是渐近自由的。

% [Tutor Note: 核对 (20.2) 的符号：$N=2$、$F=0$ 时 $b_1=2N-F=4$；$\tilde\tau\simeq-\frac{b_1}{2\pi i}\log(\mu/\Lambda)=\frac{b_1 i}{2\pi}\log(\mu/\Lambda)$，故 $\operatorname{Im}\tilde\tau=\frac{b_1}{2\pi}\log(\mu/\Lambda)>0$（当 $\mu\gg\Lambda$），与“UV 弱耦合”一致。若把符号取反，会得到负虚部，度量不正定。]

\begin{proposition}[20.1]
纯 $4d$ $\mathcal{N}=2$ $\mathfrak{su}(2)$ SYM 的经典模空间与其经典库仑分支重合，后者由条件 $[\langle\phi^\dagger\rangle,\langle\phi\rangle]=0$ 定义。每个解都规范等价于 $\langle\phi\rangle=\alpha\sigma_3$，其中 $\alpha$ 是复数。该真空期望值（VEV）把规范群自发破缺到 $\U(1)\times\Z_2$，其中 $\Z_2$ 是 $\mathfrak{su}(2)$ 的 Weyl 群。
\end{proposition}

由命题 18.25，该理论具有量子库仑分支 $\mathcal{C}$，它构成整个量子模空间。库仑分支 $\mathcal{C}$ 仍可由复参数 $\alpha$ 在规范等价意义下参数化。具体地说，剩余的离散规范对称性 $\Z_2$ 作用为 $\alpha\mapsto-\alpha$。因此：

\begin{proposition}[20.2]
$\mathcal{C}$ 上一个合适的规范不变坐标是 $u=\langle\operatorname{tr}\phi^2\rangle/2=\alpha^2$。
\end{proposition}

在低于 $\U(1)$ 下带电的有质量激发的质量的能量标度上，这些激发从理论中脱耦。例如，$W^\pm$ 玻色子由于这一自发对称性破缺而获得正比于 $\alpha$ 的质量，从而在低于 $|\alpha|$ 的能量下脱耦。因此，在极低能量下有效理论化为纯 $\mathcal{N}=2$ Maxwell 理论，它是自由理论。然而，在库仑分支的非一般点处，可以存在在低能规范群 $\U(1)$ 下带电的无质量激发；这些激发随后出现在低能有效理论中。

库仑分支 $\mathcal{C}$ 是一个特殊凯勒流形。一般地，$\mathcal{C}$ 上的低能有效理论是单个 $\mathcal{N}=2$ 矢量超场 $A$ 映入 $\mathcal{C}$ 的 $\mathcal{N}=2$ sigma 模型。其标量分量的 VEV $a$ 是 $\mathcal{C}$ 上的局部全纯坐标，即 $u$ 的局部全纯函数。概略地说，低能有效拉格朗日量的玻色部分为
\[
\mathcal{L}_{\mathrm{eff}}=\operatorname{Im}(\tau(u))\Bigl(\partial_\mu a\,\partial^\mu\bar a+F^{U(1)}_{\mu\nu}F^{\mu\nu}_{U(1)}\Bigr)+\operatorname{Re}(\tau(u))F^{U(1)}_{\mu\nu}(\star F^{\mu\nu}_{U(1)}),
\]
其中 $\tau(u)$ 是 $u\in\mathcal{C}$ 处的低能阿贝尔复化规范耦合。该拉格朗日量完全由预势 $\mathcal{F}$ 决定，后者是 $A$ 的局部全纯函数。特别地：
\[
\tau=\frac{\partial^2\mathcal{F}}{\partial a^2}=\frac{\partial a_D}{\partial a},
\]
其中 $a_D$ 是 $a$ 的磁对偶坐标，定义于式 (19.18)。

回顾 §17.3.4 末尾讨论的重要提醒：式 (20.2) 给出的 UV 规范耦合 $\tilde\tau$ 的跑动只在 $\tilde\tau$ 的虚部很大时（等价地规范耦合 $g^2$ 很小时）可靠。该条件在远大于 $|\Lambda|$ 的能量标度上满足。因此可以区分库仑分支 $\mathcal{C}$ 中两个不同的区域，如图 20.1 所示。确切地说：

% [原文 p.289]
\Fig[0.55\linewidth]{figures/fig_20_1.png}{20.1}{纯 $\mathcal{N}=2$ $\SU(2)$ SYM 库仑分支的结构。}

(I) \emph{弱耦合区域}是 $|u|\gg|\Lambda^2|$ 的区域。在该区域中，标度 $|\alpha|$ 处的规范对称性破缺发生在 UV 非阿贝尔规范动力学进入强耦合区之前。因此，该区域中 UV 理论的行为可以由从 UV 一直到标度 $|\alpha|$ 的微扰论准确描述。低于该标度时，有效理论基本变为自由理论，规范耦合的跑动可以如图 20.2 所示。

\Fig[0.40\linewidth]{figures/fig_20_2.png}{20.2}{弱耦合区域 (I) 中规范耦合的跑动。}

(II) \emph{强耦合区域}是 $|u|\lesssim|\Lambda^2|$ 的区域。在该区域中，UV 非阿贝尔规范动力学在规范对称性破缺之前就变为强耦合。因此，从 UV 一直到标度 $|\alpha|$，UV 耦合的跑动不能由式 (20.2) 的微扰表达式描述。

我们将依次分析区域 (I) 与 (II)。

% [原文 p.290]
\subsection{弱耦合区域}

在量子库仑分支 $\mathcal{C}$ 的无穷远处，量子修正变小，分支趋近经典模空间，在那里 $a=\alpha$。因此，在库仑分支的弱耦合区域中有如下关系：
\[
a=\sqrt{u}+\dots,
\]
其中省略号表示量子修正。经典地，未破缺的 $\U(1)$ 规范对称性的场强可表示为
\[
F_{\mu\nu}=\begin{pmatrix}F^{U(1)}_{\mu\nu}&0\\0&-F^{U(1)}_{\mu\nu}\end{pmatrix}=F^{U(1)}_{\mu\nu}\sigma_3,
\]
这作为我们对低能 $F^{U(1)}_{\mu\nu}$ 的定义\footnote{这里我们假定 UV 规范群是 $G=\SU(2)$。诸如 $\SO(3)_+$ 与 $\SO(3)_-$ 等其他整体变体的讨论将在本节后文给出。}。由此得到如下关系：
\[
\frac{1}{2g^2}\operatorname{tr}\bigl(F_{\mu\nu}F^{\mu\nu}\bigr)=\frac{1}{g^2}F^{U(1)}_{\mu\nu}F^{\mu\nu}_{U(1)},
\]
它表明 $g^2=4e^2$，其中 $e$ 是低能阿贝尔规范耦合。类似地，拓扑 $\Theta$ 项变为
\[
\frac{\Theta}{16\pi^2}\operatorname{tr}F_{\mu\nu}(\star F)^{\mu\nu}=\frac{\Theta}{8\pi^2}F^{U(1)}_{\mu\nu}\bigl(\star F^{U(1)}\bigr)^{\mu\nu},
\]
由此推出 $\U(1)$ 的 $\Theta$ 角为 $\Theta_{U(1)}=2\Theta$。

\begin{proposition}[20.3]
用复化耦合表示，这给出
\[
\tau:=\frac{2\pi i}{e^2}+\frac{\Theta_{U(1)}}{2\pi}=2\Bigl(\frac{4\pi i}{g^2}+\frac{\Theta}{2\pi}\Bigr)=2\tilde\tau,
\]
其中 $\tau$ 是复化低能阿贝尔规范耦合，$\tilde\tau$ 是 UV 耦合。由于 UV 规范耦合在标度 $|\alpha|$ 以下停止跑动，极红外的低能阿贝尔耦合 $\tau(\alpha)$ 满足
\begin{equation}
\tau(\alpha)\approx2\tilde\tau(\alpha)=-\frac{b_1}{\pi i}\log\Bigl(\frac{\alpha}{\Lambda}\Bigr)-\frac{2i}{\pi}\sum_k c_k\Bigl(\frac{\Lambda}{\alpha}\Bigr)^{4k},
\tag{20.3}
\end{equation}
其中当 $|\alpha|$（等价地 $|u|$）增大时非微扰瞬子修正变得可以忽略。
\end{proposition}

在弱耦合区域 (I) 中，UV 规范耦合的微扰单圈重整化是可靠的近似。把式 (20.3) 中的 $\alpha$ 替换为 $a$ 并对 $a$ 积分一次，得到 $a_D$ 的如下表达式：
% [原文 p.291]
\[
a_D=-\frac{a b_1}{\pi i}\log\Bigl(\frac{a}{\Lambda}\Bigr)+\frac{b_1a}{\pi i}+\dots.
\]

\begin{proposition}[20.4]
由于 $u\simeq a^2$，从弱耦合区域 (I) 中的一点 $u_0$ 出发、在不离开区域 (I) 的情况下绕区域 (II) 一整圈（如图 20.1 所绘），会诱导特殊坐标 $(a,a_D)$ 的如下变换：
\[
\begin{pmatrix}a\\
\]
该变换可写成矩阵形式
\[
(a\ \ a_D)\longmapsto(a\ \ a_D)\begin{pmatrix}-1&4\\0&-1\end{pmatrix}=(-a\ \ 4a-a_D).
\]
重复这一过程多次，会在 $u_0\in\mathcal{C}$ 处得到无穷多个彼此不同但物理等价的描述。重要的是，理论的所有内在特征都必须在该变换下不变。特别地，在低能 $\U(1)$ 规范群下带电的粒子的质量谱必须保持不变。
\end{proposition}

% [Tutor Note: 用 (20.4) 核对坐标变换与电荷变换的相容性：设 $(a,a_D)\mapsto(-a,-a_D+4a)$、$(n_e,n_m)\mapsto(-n_e-4n_m,-n_m)$，则 $n'_ea'+n'_ma'_D=(-n_e-4n_m)(-a)+(-n_m)(-a_D+4a)=n_ea+n_ma_D$，故 BPS 质量不变。若只变坐标不变电荷，质量会改变，所以 $M_\infty$ 对荷的作用必须由逆转置给出。]

秩 $1$ 的 $4d$ $\mathcal{N}=2$ SYM 理论在库仑分支上的 BPS 质量公式为
\begin{equation}
M=\sqrt2\,|n_ea+n_ma_D|=\sqrt2\,\Bigl|(a\ \ a_D)\begin{pmatrix}n_e\\|.
\tag{20.4}
\end{equation}
为使谱在该变换下不变，低能荷 $(n_e,n_m)$ 也必须相应地变换：
\[
\begin{pmatrix}n_e\\
\]
矩阵
\[
M_\infty=\begin{pmatrix}-1&4\\0&-1\end{pmatrix}\in\SL_2(\Z)\simeq\Sp_2(\Z)
\]
称为\emph{无穷远处的单值性}。$M_\infty\in\Sp_2(\Z)$ 这一事实保证该矩阵保持反对称的 Dirac--Schwinger--Zwanziger（DSZ）配对。

因此，即使只知道 UV 耦合的微扰单圈跑动，我们也能断定低能理论中的电磁荷不是整体地定义的，而是在绕

% [原文 p.292]
紧化库仑分支 $\mathcal{C}\simeq\C\PP^1$ 的无穷远点一圈时发生单值变换。数学上，这意味着低能荷的格构成一个带有 DSZ 配对的局部系统\footnote{拓扑空间 $X$ 上的\emph{局部系统}是局部常值层。更几何地说，它可描述为带平坦联络的丛，或基本群 $\pi_1(X)$ 的表示——这就是第 3 章讨论的 Riemann--Hilbert 对应。对格而言，局部系统是格的纤维化。}$\Gamma(u)$。

在 $\mathcal{C}$ 的每一点 $u$ 处（除有限多个奇点外），$\Gamma(u)$ 是由电磁荷 $(n_e,n_m)$ 构成的秩 $2$ 格，带有 DSZ 配对。该格在绕奇点一圈时经历非平凡的单值变换，其中一个奇点位于无穷远，其单值性为 $M_\infty$。此外，$\Gamma(u)$ 全纯地随 $u$ 变化，这由 $a$ 与 $a_D$ 的全纯性以及 BPS 质量公式所蕴含。紧化库仑分支示意于图 20.3。

\Fig[0.58\linewidth]{figures/fig_20_3.png}{20.3}{纯 $\mathcal{N}=2$ $\SU(2)$ SYM 库仑分支的结构。}

\subsection{强耦合区域}

在该区域中，人们既不能依赖式 (20.3)，也难以把库仑分支上由 UV 伴随标量 VEV 定义的自然坐标 $u$ 与自然的低能变量 $a$ 联系起来。然而，强耦合区域中发生的一切必须与我们在弱耦合区域得到的结论相容，这施加了显著的约束。

\begin{proposition}[20.5]
库仑分支具有 $\Z_2$ 对称性。
\end{proposition}

% [原文 p.293]
\begin{proof}
纯 $\mathcal{N}=2$ $\SU(2)$ SYM 具有离散 $R$ 对称性 $(\Z_8)_R$（参看 §18.3.2），在其下伴随标量 $\phi$ 的荷为 $2$。计入离散的 $\Z_2$ 规范（Weyl）对称性后，$(\Z_8)_R$ 在库仑分支上诱导出 $\Z_2$ 对称性：$u\mapsto-u$。
\end{proof}

注意 $\operatorname{Im}\tau(a)$ 必须在 $\mathcal{C}$ 上为正且单值，才能使
\begin{equation}
\dd s^2=\operatorname{Im}\tau(a)\,\dd a\,\dd\bar a
\tag{20.5}
\end{equation}
成为真正的度量。由于 $\tau$ 是 $a$ 的全纯函数，其虚部 $\operatorname{Im}\tau$ 是调和函数。此外它必须单值，否则式 (20.5) 无意义。于是调和性意味着 $\operatorname{Im}\tau$ 不能整体定义，否则它将无界。此外还有如下结论：

\begin{proposition}[20.6]
经典库仑分支 $\mathcal{C}_{\mathrm{cl}}$ 的结构（在 $\mathcal{C}$ 的原点处只有一个奇点）不可能是真正的库仑分支 $\mathcal{C}$ 的正确结构。
\end{proposition}

\begin{proof}
由 UV 规范耦合的微扰重整化我们知道，电荷格局部系统 $\Gamma(u)$ 在无穷远处有非平凡的单值性 $M_\infty$。这意味着电荷的局部系统在 $\mathcal{C}$ 中不可能平凡，因为 $\pi_1(\mathcal{C})=\pi_1(\mathcal{C}-\{\infty\})=\{\mathrm{id}\}$。在可能出现的各种拓扑复杂化之中，最简单的解或许是假定 $\mathcal{C}$ 中存在额外的奇点，$\Gamma(u)$ 绕它们发生单值变换。这个解给出一个非常有说服力的图像，强烈暗示其正确性。设 $\mathcal{C}$ 在某些点 $u_1,\dots,u_n$ 处有额外奇点，并设 $M_1,\dots,M_n\in\SL_2(\Z)$ 是绕它们一圈的单值性（取定基点 $u_0\in\mathcal{C}$ 以及以 $u_0$ 为基点、逆时针绕每个奇点 $u_i$ 的环路）。$\pi_1(\mathcal{C})$ 的平凡性意味着
\[
M_\infty^{-1}=\prod_{i=1}^{n}M_i.
\]
如果只有一个这样的额外奇点（正如 $\mathcal{C}_{\mathrm{cl}}$ 中 $W$ 玻色子变为无质量、完整的 $\SU(2)$ 规范对称性得以恢复的情形），那么由命题 20.5 它必在原点。于是低能耦合 $\tau$ 作为 VEV $a$ 的函数将具有形式
\[
\tau(a)=-\frac{1}{\pi i}\log\frac{a}{\Lambda}+f(a),
\]
其中 $f$ 是 $a$ 的亚纯函数、在原点有单极点。这又意味着 $\operatorname{Im}\tau$ 无界，从而不是正定的。
\end{proof}

因此 $\mathcal{C}$ 最简单的奇点结构由位于相反取值 $\pm u_1$ 的两个奇点组成。由于 $\Lambda$ 是理论中唯一的标度，

% [原文 p.294]
可以合理期望 $u_1\propto\Lambda^2$。事实证明，假定在 $\pm2\Lambda^2$ 处存在两个奇点会得到一个非常有说服力的框架。$\mathcal{C}$ 的建议结构示于图 20.4，其中 $M_1$ 与 $M_2$ 已改记为 $M_\pm$。

\Fig[0.55\linewidth]{figures/fig_20_4.png}{20.4}{纯 $\mathcal{N}=2$ $\SU(2)$ SYM 的库仑分支。}

设 $M_+$ 与 $M_-$ 分别是 $(a,a_D)$（等价地电荷局部系统）在 $2\Lambda^2$ 与 $-2\Lambda^2$ 处单值性的矩阵。这些矩阵必须满足
\[
M_\infty=M_+M_-.
\]
此外，由于 $\Z_2$ 对称性把 $2\Lambda^2$ 与 $-2\Lambda^2$ 互换，$M_+$ 与 $M_-$ 必须在 $\SL_2(\Z)$ 中共轭。可以直接验证：

\begin{proposition}[20.7]
矩阵
\begin{equation}
M_+=\begin{pmatrix}1&0\\-1&1\end{pmatrix}\qquad\text{与}\qquad M_-=\begin{pmatrix}-1&4\\-1&3\end{pmatrix}
\tag{20.6}
\end{equation}
同时满足这两个约束。
\end{proposition}

% [Tutor Note: 直接相乘核对 $M_\infty=M_+M_-$：$\begin{pmatrix}1&0\\-1&1\end{pmatrix}\begin{pmatrix}-1&4\\-1&3\end{pmatrix}=\begin{pmatrix}-1&4\\0&-1\end{pmatrix}$。又 $M_\pm$ 的行列式均为 $1$、迹均为 $2$，故二者在 $\SL_2(\Z)$ 中共轭，与 $\Z_2$ 对称性 $u\mapsto-u$ 互换 $2\Lambda^2$ 与 $-2\Lambda^2$ 相容。]

对超对称真空模空间中的奇点，一个自然的解释是：它们标志着由于出现额外无质量粒子而导致的低能描述失效，如第 17 章在 $\mathcal{N}=1$ 理论语境中所讨论的那样。例如，经典库仑分支 $\mathcal{C}_{\mathrm{cl}}$ 原点的奇点出现，是因为 $W^\pm$ 玻色子在该点变为无质量。然而，所建议的奇点位于非零 VEV $\langle\operatorname{tr}\phi^2\rangle\propto\Lambda^2$ 处的图景，与这些奇点联系于无质量规范玻色子的想法相冲突 [313]。相反，它暗示量子库仑分支 $\mathcal{C}$ 上的两个奇点源于 BPS 超多重态变为无质量。

设 $M_{(n_e,n_m)}\in\SL_2(\Z)$ 是由电磁荷为 $(n_e,n_m)$ 的低能 BPS 超多重态变为无质量所诱导的单值性。该单值性必须使 $(n_e,n_m)$ 保持不变。

% [原文 p.295]
BPS 质量公式 $M=\sqrt2|n_ea+n_ma_D|$ 表明，磁单极子在满足 $a_D=0$ 的点 $u_s\in\mathcal{C}$ 处变为无质量。由于磁单极子以非局域方式与光子场 $A$ 耦合，不可能存在同时含有 $A$ 与单极子场的低能有效拉格朗日量。然而，单极子与 $A_D$ 局部耦合。因此可以作形如式 (19.20) 的对偶变换，用磁对偶场 $A_D$ 与单极子场来表达低能拉格朗日量。

注意当逼近 $u_s$ 时低能耦合 $e_D$ 趋于零，而 $e$ 在彼处发散。等价地，当 $a_D\to0$ 时 $\operatorname{Im}\tau_D\to\infty$。因此 $u_s$ 附近的低能理论是 $\mathcal{N}=2$ Maxwell 理论，但规范场为 $A_D$、轻物质场为单极子。这表明在 $\mathcal{C}$ 上接近 $u_s$ 处，$a_D$ 比 $a$ 更适合作为局部坐标。此外，在 $u_s$ 附近可写
\[
\tau_D\simeq-\frac{i}{\pi}\ln a_D,
\]
其中
\[
a_D=c_0(u-u_s).
\]
由于 $\tau_D=-\dd a/\dd a_D$，这给出 $a(u)$ 的如下表达式：
\[
a(u)\simeq a_0+\frac{i}{2\pi}c_0(u-u_s)\ln(u-u_s),
\]
其中常数 $a_0$ 不能为零，否则电粒子也会在 $u_s$ 处无质量，从而使以 $A_D$ 与单极子表述的低能拉格朗日量描述失效。当 $u$ 逆时针绕 $u_s$ 一圈时，特殊坐标 $(a,a_D)$ 经历单值变换
\[
(a\ \ a_D)\longmapsto(a-a_D\ \ a_D)=(a\ \ a_D)M_+,
\]
其中 $M_+$ 由式 (20.6) 定义。于是：

\begin{proposition}[20.8]
单值性 $M_+$ 由在 $2\Lambda^2$ 处变为无质量的单极子所诱导。
\end{proposition}

现在考虑一个荷为 $(2,1)$ 的低能 BPS 双子。其质量在满足 $2a+a_D=0$ 的点 $u_s\in\mathcal{C}$ 处为零。在这样的点附近，$2a+a_D$ 是 $\mathcal{C}$ 上的好坐标：
\[
2a+a_D=c'_0(u-u_s),
\]
且当 $2a+a_D\to0$ 时 $\operatorname{Im}\tau_-=\operatorname{Im}(2\tau+1)/(\tau+1)\to\infty$。由于 $\tau_-=\dd(a+a_D)/\dd(2a+a_D)$，可以推出
\[
(a+a_D)(u)\simeq b_0-\frac{i}{2\pi}c'_0(u-u_s)\ln(u-u_s).
\]
于是当逆时针绕 $u_s$ 一圈时，特殊坐标 $(a,a_D)$ 经历单值变换

% [原文 p.296]
\[
(a\ \ a_D)\longmapsto(a\ \ a_D)M_-=(-a-a_D\ \ 4a+3a_D),
\]
其中 $M_-$ 由式 (20.6) 定义。于是：

\begin{proposition}[20.9]
单值性 $M_-$ 由荷为 $(2,1)$ 的双子变为无质量所诱导。
\end{proposition}

推广上述推理可得：

\begin{proposition}[20.10]
由荷为 $(n_e,n_m)$ 的双子在 $\mathcal{C}$ 中某点变为无质量所诱导的 $(a,a_D)$ 的单值性为
\[
M_{(n_e,n_m)}=\begin{pmatrix}1-n_en_m&n_e^2\\-n_m^2&1+n_en_m\end{pmatrix}.
\]
\end{proposition}

\begin{remark}[20.11]
注意单值性依赖于以 $u_0$ 为基点、绕各奇点的环路的选取。例如，若在绕 $2\Lambda^2$ 处奇点之前逆时针绕强耦合区域 (II)、之后顺时针绕它，则单值性表现为 $M_\infty^{-1}M_+M_\infty$ 而不是 $M_+$。由于
\[
M_\infty^{-1}M_+M_\infty=\begin{pmatrix}-3&16\\-1&5\end{pmatrix},
\]
$2\Lambda^2$ 处的奇点现在似乎对应于荷为 $(4,1)$ 的双子变为无质量，而不是磁单极子。归根结底，比具体的单值性 $M_\infty,M_\pm$ 更重要的是平坦 $\SL_2(\Z)$ 丛的整体结构，而 $(a,a_D)$ 是该丛的一个截面。
\end{remark}

\subsection{从 $\mathcal{N}=2$ 破缺到 $\mathcal{N}=1$}

对 $\mathcal{N}=2$ 矢量超场中的手征超场添加质量项 $m\operatorname{tr}\Phi^2$ 会破坏纯 $\mathcal{N}=2$ $\SU(2)$ SYM 的两个超对称中的一个，使低能理论化为纯 $\mathcal{N}=1$ $\mathfrak{su}(2)$ SYM。如 §17.3.5 所讨论，纯 $\mathcal{N}=1$ $\SU(2)$ SYM 具有手征 $\Z_4$ $R$ 对称性，并被认为会禁闭，产生质量能隙，并显示手征 $\Z_4$ $R$ 对称性自发破缺到 $\Z_2$，从而给出两个不同的真空。

若采用纯 $\mathcal{N}=2$ $\mathfrak{su}(2)$ SYM 的库仑分支 $\mathcal{C}$ 在 $\pm2\Lambda^2$ 处具有两个奇点这一猜想，我们就能理解：引入保持 $\mathcal{N}=1$ 的质量项如何提升模空间 $\mathcal{C}$，却保留对应于额外无质量自由度的两个临界点 [313]。在这些点处，低能规范场 $A$ 因希格斯机制获得质量，理论产生质量能隙。此外，纯 $\mathcal{N}=2$

% [原文 p.297]
$\SU(2)$ 的 $\Z_4$ $R$ 对称性确实自发破缺到 $\Z_2$，因为在 $\mathcal{C}$ 上非平凡作用的 $\Z_2$ 子群交换 $2\Lambda^2$ 与 $-2\Lambda^2$。

这一推理为量子库仑分支 $\mathcal{C}$ 中所猜想的奇点结构提供了有力的验证。

\subsection{Seiberg–Witten 曲线}

我们希望构造 $\mathcal{C}$ 上的全纯函数 $a$ 与 $a_D$，使 $(a,a_D)$ 在 $\pm u_0$ 处经历单值性 $M_\pm$、在无穷远处经历 $M_\infty$。这是 $\mathcal{C}$ 上关于平坦 $\SL_2(\Z)$ 丛的一个 Riemann--Hilbert 问题。由于渐近条件由 $\mathcal{C}$ 弱耦合区域的半经典分析固定，解是唯一的。事实证明，$a$ 与 $a_D$ 可以用周期积分整洁地解释 [313]。

\begin{proposition}[20.12]
单值性 $M_\pm$ 与 $M_\infty$ 生成 $\mathrm{PSL}_2(\Z)$ 的子群 $\Gamma_0(4)$：
\[
\Gamma_0(4)=\Bigl\{\begin{pmatrix}a&b\\c&d\end{pmatrix}\in\SL_2(\Z)\ \Big|\ b\equiv0\ [4]\Bigr\},
\]
它分类带有一个 $4$ 阶循环子群的复环面。上半平面 $\HH$ 对 $\Gamma_0(4)$ 的商\footnote{该商称为\emph{经典模曲线}，记作 $X_0(4)$。}具有与 $\mathcal{C}$ 所建议的相同的奇点结构。
\end{proposition}

这提示用 $\mathcal{C}$ 上的复环面族来构造 $a$ 与 $a_D$，该族由模空间 $\HH/\Gamma_0(4)$ 上的重言纤维化给出。考虑如下方程：
\begin{equation}
(\Sigma_u):\quad \frac{\Lambda^2}{z}+\Lambda^2 z=x^2-u,
\tag{20.7}
\end{equation}
其中 $z\in\C^\times$、$x\in\C$。这在 $\C^\times_z\times\C_x\times\C_u$ 中定义了一个复曲面，可以通过齐次化方程来紧化，特别允许 $z\in\C\PP^1$。该族具有明显的 $\Z_4$ 对称性，在其下 $(u,x,z)\mapsto(-u,ix,-z)$，它投影为 $\mathcal{C}$ 上的 $\Z_2$ 对称性 $u\mapsto-u$。

对一切一般的 $u$，式 (20.7) 定义了椭圆曲线 $\Sigma_u\subset\C\PP^1_z\times\C_x$，它是 $\C\PP^1_z$ 的 $2$ 次分歧覆叠，在 $0$、$\infty$ 以及 $z_\pm$ 处有四个 $2$ 阶分歧点，其中
\begin{equation}
z_\pm=\frac{-u\pm\sqrt{u^2-4\Lambda^4}}{2\Lambda^2}.
\tag{20.8}
\end{equation}
回顾库仑分支上的度量为
\[
\dd s^2=\operatorname{Im}(\tau)\,\dd u\,\dd\bar u,
\]
其中

% [原文 p.298]
\begin{equation}
\tau(u)=\frac{\dd a_D}{\dd a}=\frac{\dd a_D/\dd u}{\dd a/\dd u}
\tag{20.9}
\end{equation}
必须有正定的虚部。

有一种典范的方式从式 (20.7) 定义的椭圆曲线族在 $\mathcal{C}$ 上构造这样的 $\tau(u)$：对每个 $u$，取 $\tau(u)\in\HH$ 为 $\Sigma_u$ 的复模。具体地说，每个 $\Sigma_u$ 容许唯一的（相差常数倍）全纯微分 $\omega_u=\dd z/xz$。设 $A$ 与 $B$ 是 $H_1(\Sigma_u,\Z)$ 的辛基（即要求 $A\cdot B=1$，其中 $\cdot$ 是相交配对），并令
\[
\omega_A=\oint_A\omega,\qquad \omega_B=\oint_B\omega.
\]
比值 $\omega_B/\omega_A$ 是 $\Sigma_u$ 的复模，它位于 $\HH$ 中因而具有正的虚部。这是 Riemann 双线性关系的一个特例。

\begin{definition}[20.13]
设 $\lambda$ 是一族亚纯微分，其中 $\lambda_u$ 是 $\Sigma_u$ 上留数为零的亚纯微分，且满足 $\dd\lambda/\dd u=\omega_u$。于是令
\[
a=\oint_A\lambda,\qquad a_D=\oint_B\lambda,
\]
就可以如式 (20.9) 那样表达 $\tau(u)$。
\end{definition}

因此式 (20.9) 的周期是 $\mathcal{C}$ 上特殊坐标 $a$ 与 $a_D$ 的有效候选。事实上，$\lambda$ 在模去恰当形式的范围内由 $a$ 与 $a_D$ 的渐近行为完全确定：
\[
\lambda=x\frac{\dd z}{z}.
\]
该微分满足 $\dd\lambda/\dd u=\omega_u$。关于式 (20.9) 定义的周期 $a$ 与 $a_D$ 确实可以等同于 $\mathcal{C}$ 上的特殊坐标 $a$ 与 $a_D$ 的更明确验证，参见 [313, 314, 339]；以下我们假定这一等同成立。

\begin{remark}[20.14]
这一等同特别意味着电性带电粒子从不在 $\mathcal{C}$ 上变为无质量，因为可以算出 $a$ 从不为零，尤其在 $u=0$ 处不为零。因此，$\mathcal{C}_{\mathrm{cl}}$ 原点的经典奇点在量子库仑分支 $\mathcal{C}$ 中实际上已不复存在，这与上文对强耦合区域的分析及其猜想的奇点结构一致。
\end{remark}

式 (20.7) 的纤维化清楚地表明 $\pm2\Lambda^2$ 是奇点，因为式 (20.8) 的分歧点 $z_+$ 与 $z_-$ 在那里相撞。在这些奇点处，$1$-圈 $A$ 与 $B$（从而 $a$ 与 $a_D$）的单值性由 Picard--Lefschetz 单值性公式给出：在一个 $1$-圈 $\nu\in H_1(\Sigma_u,\Z)$ 消失的点处，单值性在 $H_1(\Sigma_u,\Z)$ 上的作用为
\[
\gamma\longmapsto\gamma-(\gamma\cdot\nu)\nu,\qquad\text{对一切 }\gamma\in H_1(\Sigma_u,\Z).
\]
取定与 $a,a_D$ 在无穷远处的渐近行为相容的 $1$-圈 $A,B$ 后，容易验证 $2\Lambda^2$（相应地 $-2\Lambda^2$）处的消失圈是 $B$（相应地 $2A+B$），从而得到式 (20.6) 的矩阵 $M_+$ 与 $M_-$。

% [原文 p.299]
\subsection{BPS 谱与穿墙}

在任一点 $u\in\mathcal{C}$，低能电磁荷为 $(n_e,n_m)$ 的 BPS 态的中心荷由重整化群不变的表达式给出：
\[
\mathcal{Z}=n_ea+n_ma_D,
\]
从而 BPS 质量公式为
\[
|M|=\sqrt2\,|n_ea+n_ma_D|.
\]
Seiberg--Witten 解允许计算任何假想的、带电磁荷 $(n_e,n_m)$ 的低能 BPS 粒子的中心荷，从而计算其质量。然而，虽然中心荷可以对潜在的 BPS 态计算，这并不保证它们存在于低能理论中。与 $\mathcal{N}=4$ SYM 理论中人们期望对所有可能的低能荷都存在 BPS 态相反，在纯 $4d$ $\mathcal{N}=2$ $\SU(2)$ SYM 中，BPS 谱受到更多限制。

在弱耦合区域 (I) 中，BPS 谱可以用 't Hooft--Polyakov 与 Julia--Zee 的方法半经典地计算。它由荷为 $(\pm2,0)$ 的 $W^\pm$ 玻色子、荷为 $(0,\pm1)$ 的磁单极子以及荷为 $\pm(2k,\pm1)$（$k\in\N$）的双子组成。$W^\pm$ 玻色子属于 BPS 矢量多重态，而其他所有单极子与双子都处于 BPS 超多重态中。

\begin{proposition}[20.15]
半经典 BPS 谱在单值性 $M_\infty$\footnote{然而单个粒子并不如此。例如，在图 20.1 所描绘的过程中，一个磁单极子变为荷为 $(-4,-1)$ 的双子，而荷为 $(2,-1)$ 的双子变为荷为 $(2,1)$ 的双子。}下整体不变。然而它在 $M_\pm$ 下不变，因此在 $\Gamma_0(4)$ 下也不变。
\end{proposition}

例如，在 $M_+$ 下，荷为 $(2,1)$ 的双子变为荷为 $(2,3)$ 的双子，而后者在弱耦合区域中不存在。这显然对在 $\mathcal{N}=2$ SYM 理论中实现 S 对偶构成挑战。然而这一结论为时过早，因为库仑分支 $\mathcal{C}$ 内会发生穿墙，使得 BPS 谱可以在跨越不同区域时改变。

回顾命题 19.34：三角不等式
\[
|\mathcal{Z}_1+\mathcal{Z}_2|\le|\mathcal{Z}_1|+|\mathcal{Z}_2|
\]
意味着荷为 $(n_e,n_m)$（其中 $n_e$ 与 $n_m$ 互素）的 BPS 态一般对衰变为两个或更多粒子是稳定的。这是因为一般地 $a/a_D\notin\R$，从而不等式严格。然而，当 $a/a_D\in\R$ 时不等式取等，BPS 态只是边缘稳定的。在这种情形下衰变可能发生，因为 BPS 态可以分裂为多个粒子而仍保持 BPS 条件。这一穿墙现象首先是在

% [原文 p.300]
二维量子场论的语境中被发现的 [74, 75]。Seiberg 与 Witten 的工作后来表明，穿墙也发生在 $4d$ $\mathcal{N}=2$ 理论中。

\begin{definition}[20.16]
条件 $a/a_D\in\R$ 定义了 $\mathcal{C}$ 中实余维 $1$ 的子流形，称为\emph{边缘稳定墙}。
\end{definition}

在这些墙的两侧，BPS 态的谱保持不变，尽管某个特定 BPS 态的中心荷（从而其质量）可以随 $\mathcal{C}$ 中的点而变。当跨越一道边缘稳定墙时，可能发生衰变，导致 BPS 谱的跳变。在纯 $\mathcal{N}=2$ $\SU(2)$ SYM 中，存在这样一道墙：它是经过 $\pm2\Lambda^2$ 的一条曲线，如图 20.5 所示。

\Fig[0.72\linewidth]{figures/fig_20_5.png}{20.5}{纯 $\mathcal{N}=2$ $\SU(2)$ SYM 库仑分支上的边缘稳定墙。}

无质量 BPS 态只能存在于 $\mathcal{C}$ 上，因为 $n_ea+n_ma_D=0$ 意味着 $a/a_D\in\Q\subset\R$。此外可以证明，当沿该墙从 $-2\Lambda^2$ 出发、先在上半平面、再在下半平面经过 $2\Lambda^2$ 回到自身时，$a_D/a$ 从 $-2$ 单调地变到 $2$。这意味着若 $n_e/n_m\in[-2,2]$，则荷为 $(n_e,n_m)$ 的 BPS 态必在曲线上的某处变为无质量。若一个 BPS 态在墙上变为无质量，则它在该墙两侧都存在，因为它是跨越该墙时出现的唯一无质量粒子，故它是稳定的。因此，既然我们知道荷为 $\pm(2,\pm1)$ 的单极子与双子在 $\mathcal{C}$ 上变为无质量，它们就是弱耦合区域中以及边缘稳定墙内部的稳定 BPS 态 [127]。

\begin{proposition}[20.17]
在边缘稳定曲线之外，BPS 谱由 $W^\pm$ 玻色子、磁单极子以及所有荷为 $\pm(2k,\pm1)$（$k\in\Z$）的双子组成。在曲线之内，BPS 谱只由磁单极子与特定的双子组成 [127]。可以取实轴上从 $-\infty$ 到 $1$ 的线段为 $M_+$ 的分支割线。采用图 20.5 的记号，谱如下：
\begin{itemize}
\item $R_W$：$W^\pm$ 玻色子 $\pm(2,0)$、磁单极子 $\pm(0,1)$、双子 $\pm(2k,\pm1)$（$k\in\Z$），
\item $R^+_S$：磁单极子 $\pm(0,1)$、双子 $\pm(2,-1)$，
\item $R^-_S$：磁单极子 $\pm(0,1)$、双子 $\pm(2,1)$。
\end{itemize}
\end{proposition}

% [原文 p.301]
$R_W$ 与 $R_S$ 中 BPS 谱的差异表明穿墙确实发生。例如，考虑如图 20.5 中位于边缘稳定墙之外一点的荷为 $(2,1)$ 的双子。当逆时针绕 $2\Lambda^2$ 一圈并首次穿过该墙时，双子 $(2,1)$ 不再存在于 $R^+_S$ 中。因此它必须衰变：
\[
(2,1)\longmapsto(2,-1)+2\cdot(0,1).
\]
由于线段 $]-\infty,-1]$ 表示 $M_+$ 的分支割线，穿过它导致
\[
(2,-1)+2\cdot(0,1)\longmapsto(2,1)+2\cdot(0,1).
\]
因此，即使在 $R_S$ 中不存在荷为 $(2,3)$ 的 BPS 双子，穿墙也意味着单个 BPS 粒子在 $M_\pm$ 下可以变成一对 BPS 粒子。于是，尽管比 $4d$ $\mathcal{N}=4$ SYM 中更微妙，弱耦合 BPS 谱的具体形式并不与在 $4d$ $\mathcal{N}=2$ 纯 $\mathfrak{su}(2)$ SYM 中实现 S 对偶相矛盾。

\subsection{整体变体}

在纯 $4d$ $\mathcal{N}=2$ $\mathfrak{su}(2)$ SYM 理论中，物质内容完全位于规范群的伴随表示中，而式 (20.1) 的拉格朗日量对 $G=\SU(2)$ 与 $G=\SO(3)=\SU(2)/\Z_2$ 完全相同。此外 $\SO(3)$ 理论不止一个，而是有两个不同的理论，记作 $\SO(3)_+$ 与 $\SO(3)_-$，它们由线算子内容的不同而区分。这些称为\emph{整体变体}，如定义 19.26 所定义。

\begin{proposition}[20.18]
纯 $4d$ $\mathcal{N}=2$ $\mathfrak{su}(2)$ SYM 的这些整体变体之间的区别反映在 $a$ 与 $a_D$ 的定义中，从而反映在 Seiberg--Witten 曲线上 $A$ 与 $B$ 圈的选取中。这保证了低能电磁荷对应于低能 $\U(1)$ 到 UV 规范群的适当嵌入 [339, Sect. 4.5]。
\end{proposition}

当 $G=\SU(2)$ 时，电粒子（或 Wilson 线）可以处于 $\SU(2)$ 的任何表示中，从而得到关于低能 $\U(1)$ 的任何整数值电荷。$W^\pm$ 玻色子的低能荷为 $\pm(2,0)$ 这一事实表明它们的电荷是最小允许值的两倍，这与它们在 $\SU(2)$ 伴随表示下变换（该表示具有 $N$-ality $2$）相一致。此情形下的磁单极子（或 't Hooft 线）对应于 $\SU(2)$ 的余权格，基本单极子的低能磁荷为 $1$，即最小允许值。

当 $G=\SO(3)$ 时，电粒子（或 Wilson 线）必须在 $\SU(2)$ 的具有 $N$-ality $2$ 的表示下变换。低能荷的这种重新定义保证电荷 $1$ 是对应于理论中物理上可实现的粒子（或 Wilson 线）的最小电荷。相应地，$W^\pm$ 玻色子现在

% [原文 p.302]
带荷 $\pm(1,0)$。$G=\SO(3)$ 的磁单极子（或 't Hooft 线）对应于 $\SO(3)$ 的余权格，它与 $\SU(2)$ 的权格对偶。此情形下在 $2\Lambda^2$ 处凝聚的单极子具有最小允许磁荷的两倍，从而低能荷为 $\pm(0,2)$。具体地（详见 [90]）：
\begin{itemize}
\item 整体变体 $\SO(3)_+$ 通过重新定义变量 $(a',a'_D)=(2a,a_D/2)$（等价地，重新定义圈 $(A',B')=(2A,B/2)$）实现。这导致低能荷的重新定义 $(n'_e,n'_m)=(n_e/2,2n_m)$。在这个新框架中，在 $2\Lambda^2$ 处变为无质量的单极子荷为 $(0,2)$，相应的单值性矩阵 $M_+$ 为
\[
M^{\SO(3)_+}_+=\begin{pmatrix}1&0\\-4&1\end{pmatrix}.
\]
类似地，在此重新定义下有
\[
M^{\SO(3)_+}_-=\begin{pmatrix}-1&1\\-4&3\end{pmatrix},\qquad
M^{\SO(3)_+}_\infty=\begin{pmatrix}-1&1\\0&-1\end{pmatrix}.
\]
\item 整体变体 $\SO(3)_-$ 通过重新定义变量 $(a'',a''_D)=(2a,a_D/2+a)$（等价地，重新定义圈 $(A'',B'')=(2A,B/2+A)$）实现。这对应于低能荷的重新定义 $(n''_e,n''_m)=(n_e/2-n_m,2n_m)$。单值性矩阵变为
\[
M^{\SO(3)_-}_+=\begin{pmatrix}3&1\\-4&-1\end{pmatrix},\qquad
M^{\SO(3)_-}_-=\begin{pmatrix}1&0\\-4&1\end{pmatrix},\qquad
M^{\SO(3)_-}_\infty=\begin{pmatrix}-1&1\\0&-1\end{pmatrix}.
\]
\end{itemize}

% [Tutor Note: 核对整体变体的荷重定义：令 $(a',a'_D)=(2a,a_D/2)$、$(n'_e,n'_m)=(n_e/2,2n_m)$，则 $n'_ea'+n'_ma'_D=n_ea+n_ma_D$，BPS 质量公式形式不变；$\SO(3)_-$ 用 $(a'',a''_D)=(2a,a_D/2+a)$ 与 $(n''_e,n''_m)=(n_e/2-n_m,2n_m)$ 同样保持 $n''_ea''+n''_ma''_D=n_ea+n_ma_D$。]

\section{带基本物质的 $\SU(2)$ 超 Yang–Mills 理论}

拉格朗日 $4d$ $\mathcal{N}=2$ 规范理论中复化规范耦合的重整化在单圈精确，而带 $F>0$ 个基本超多重态的 $4d$ $\mathcal{N}=2$ $\SU(2)$ SYM\footnote{这里理论的整体形式没有歧义，因为 $\SU(2)$ 的基本表示不是 $\SO(3)$ 的表示。}的单圈系数为 $b_1=4-F$（参看例 18.19）。因此，对 $1\le F\le3$，理论是渐近自由的。当 $F=4$ 时，规范耦合的微扰重整化变为零。这暗示带四个（无质量）基本超多重态的四维 $\mathcal{N}=2$ $\SU(2)$ SYM 理论可能是共形不变的——这一假设有大量证据支持。对 $F>4$，理论是红外自由的，不能作为 UV 理论良好定义。

% [原文 p.303]
因此我们关注带 $1\le F\le4$ 个基本超多重态的 $4d$ $\mathcal{N}=2$ $\SU(2)$ SYM 理论，并对 $F\le3$ 用 $\Lambda_F$ 表示理论的动力学标度。

用 $\mathcal{N}=1$ 语言，每个超多重态由两个手征超场 $Q^a_i$ 与 $\tilde Q^i_a$ 组成，其中 $a=1,2$ 是色指标，$i=1,\dots,F$ 是味指标。这些超多重态可以有质量，其裸质量项 $m^i_j$ 在味群的伴随表示中变换。$\SU(2)_R$ 下的不变性施加如式 (18.6) 的条件 $[m,m^\dagger]=0$，它又意味着质量矩阵可对角化。因此不失一般性，超势可写为
\begin{equation}
W=\sqrt2\,\tilde Q^a_i\Phi^b_aQ^i_b+\sum_i m_i\tilde Q^a_iQ^i_a,\qquad m_i\in\C,
\tag{20.10}
\end{equation}
其中 $\Phi$ 是 $\mathcal{N}=2$ 矢量超场中的 $\mathcal{N}=1$ 伴随手征超场。此外 $[m,m^\dagger]=0$ 意味着一般的质量矩阵 $m$ 把味对称群破缺到其 Cartan 子群。

\subsection{整体对称性}

设所有超多重态都无质量。则：

\begin{proposition}[20.19]
理论的经典味对称群是 $\mathrm{O}(2F)$。
\end{proposition}

\begin{proof}
这来自如下事实：规范群 $\SU(2)$ 的基本表示（场 $Q^a_i$ 在其下变换）同构于场 $\tilde Q^i_a$ 所变换的反基本表示。具体地说，$\SU(2)$ 的基本表示是四元数的。更一般地，在任何规范群 $G$ 的四元数表示中变换的 $F$ 个无质量超多重态导致 $\mathrm{O}(2F)$ 味对称性。类似地，在实表示中变换的超多重态给出辛味对称群。此外，味对称性是 $\mathrm{O}(2F)$ 而非 $\SO(2F)$，是因为存在交换 $Q_1$ 与 $\tilde Q^1$ 的 $\Z_2$ 因子 $\rho$。
\end{proof}

\begin{remark}[20.20]
当质量矩阵 $m$ 非零时，经典味对称群约化为 $\mathrm{O}(2F)$ 的子群。在质量矩阵一般的极端情形下，对称群进一步破缺到 $\mathrm{O}(2F)$ 的 Cartan 子群。
\end{remark}

经典理论的整体对称群由 $\mathrm{O}(2F)\times\SU(2)_R\times\U(1)_R$ 与洛伦兹群之积的一个商给出 [314, Sect. 3]。出现这个商，是因为元素 $-1=(-1)^F\in\U(1)_R$ 实际上属于洛伦兹群，且该元素 $-1\in\U(1)_R$ 与 $\SU(2)_R$ 中心的生成元的组合必须与味群 $\mathrm{O}(2F)$ 的 $\Z_2$ 中心生成元等同。

所有带 $1\le F\le4$ 个基本超多重态的四维 $\mathcal{N}=2$ $\SU(2)$ 理论都容许复一维经典库仑分支，在其上规范群

% [原文 p.304]
自发破缺到 $\U(1)$。由命题 18.25，这些理论还容许量子库仑分支 $\mathcal{C}^{(F)}$，它是特殊凯勒流形。如 §20.1，参数 $u=\langle\operatorname{tr}\phi^2\rangle/2$ 作为 $\mathcal{C}^{(F)}$ 上的自然坐标。由于带 $1\le F\le3$ 个基本超多重态的四维 $\mathcal{N}=2$ $\SU(2)$ SYM 是渐近自由的，人们再次在 $\mathcal{C}^{(F)}$ 中区分弱耦合区域与强耦合区域。

当 $F\ge2$ 且至少两个超多重态具有相同质量时，理论还显示一个希格斯分支——一个超凯勒流形。希格斯分支总是经典精确的。对 $F=2$ 个无质量超多重态，希格斯分支由两个 $\C^2/\Z_2$ 分支组成，它们彼此相交，并在其顶点与 $\mathcal{C}^{(2)}$ 的原点相交。对 $F\ge3$ 个无质量超多重态，希格斯分支由单一分支组成，它与库仑分支 $\mathcal{C}^{(F)}$ 在原点相交。由于我们的主要兴趣在于发生在库仑分支上的 BPS 态与穿墙现象，我们不再探讨希格斯分支，而参见 [314, Sect. 3] 与 [339, Chap. 7]。

\begin{proposition}[20.21]
（[314, Sect. 4]）量子力学上，当 $1\le F\le3$ 时，经典整体对称群中的 $(\Z_2)_\rho\times\U(1)_R$ 因子反常：只有 $\Z_{4(4-F)}$ 子群作为量子理论的真正对称性存活下来。其生成元作用为
\[
\Phi\mapsto\exp\Bigl(\frac{i\pi}{4-F}\Bigr)\Phi,\qquad
Q_i\mapsto\exp\Bigl(-\frac{i\pi}{2(4-F)}\Bigr)Q_i,\qquad
\tilde Q_i\mapsto\exp\Bigl(-\frac{i\pi}{2(4-F)}\Bigr)\tilde Q_i.
\]
注意由于 $\Phi$ 在 $\Z_{4(4-F)}$ 下的荷为 $2$，$u$ 的荷为 $4$；因此 $\Z_{4(4-F)}$ 在 $\mathcal{C}^{(F)}$ 的一般点处自发破缺到 $\Z_4$，且 $\mathcal{C}^{(F)}$ 具有 $\Z_{4-F}$ 对称性。
\end{proposition}

\subsection{中心荷与单值性}

回顾 $\mathcal{N}=2$ 超庞加莱代数的中心荷 $\mathcal{Z}$ 是阿贝尔守恒荷的线性组合。在库仑分支 $\mathcal{C}^{(F)}$ 上，低能电磁荷 $(n_e,n_m)$ 如前面所讨论对 $\mathcal{Z}$ 有贡献。此外，其他阿贝尔守恒荷——例如来自有质量超多重态的那些——也对 $\mathcal{Z}$ 有贡献。

例如，在带一个裸质量 $m_1\gg\Lambda_1$ 的有质量超多重态的四维 $\mathcal{N}=2$ $\SU(2)$ SYM 中，在库仑分支 $\mathcal{C}^{(F)}$ 的弱耦合区域（即 $u\gg\Lambda_1^2$），超多重态在 $a=\pm m_1/\sqrt2$ 时变为无质量，参看式 (18.18)。这个 BPS 态具有低能电磁荷 $(n_e,n_m)=\pm(1,0)$ 与味荷 $\pm1$，因此可知
\begin{equation}
\mathcal{Z}=n_ea+n_ma_D+S\frac{m_1}{\sqrt2},
\tag{20.11}
\end{equation}
其中 $S$ 是味荷。在 $\mathcal{C}^{(F)}$ 中 $a=\pm m_1/\sqrt2$ 附近的点处有

% [原文 p.305]
\[
a\simeq\frac{m_1}{\sqrt2},\qquad
a_D\simeq c-\frac{i}{2\pi}\Bigl(a-\frac{m_1}{\sqrt2}\Bigr)\ln\Bigl(a-\frac{m_1}{\sqrt2}\Bigr),
\]
因为轻超多重态对数地贡献 $a_D$。因此单值性为
\begin{equation}
a\mapsto a,\qquad a_D\mapsto a_D+a-\frac{m_1}{\sqrt2}.
\tag{20.12}
\end{equation}
特殊坐标 $(a,a_D)$ 在单值变换下并不只是按 $\SL_2(\Z)$ 变换：它们还会获得可加常数。这与保持度量形式一般形状的对偶变换的一般结构一致，如式 (19.20) 所讨论。用矩阵形式：
\[
\Bigl(a\ \ a_D\ \ \frac{m_1}{\sqrt2}\Bigr)\longmapsto
\Bigl(a\ \ a_D\ \ \frac{m_1}{\sqrt2}\Bigr)
\begin{pmatrix}1&1&0\\0&1&0\\0&-1&1\end{pmatrix},
\]
其中最后一列的结构保证 $m_1/\sqrt2$ 在单值变换下不变（$m_1$ 是理论中非动力学的参数）。对荷的作用由逆矩阵给出：
\[
\begin{pmatrix}n_e\\
\begin{pmatrix}1&-1&0\\0&1&0\\0&1&1\end{pmatrix}
\begin{pmatrix}n_e\\
\begin{pmatrix}n_e-n_m\\
\]
裸质量 $m_i$ 在单值变换下不变，意味着规范荷可以贡献味荷，而味荷从不贡献规范荷。

% [Tutor Note: 核对三联单值变换下中心荷 (20.11) 不变：$(a,a_D,m_1/\sqrt2)\mapsto(a,\ a_D+a-m_1/\sqrt2,\ m_1/\sqrt2)$，而 $(n_e,n_m,S)\mapsto(n_e-n_m,\ n_m,\ S+n_m)$，于是 $\mathcal{Z}'=(n_e-n_m)a+n_m(a_D+a-m_1/\sqrt2)+(S+n_m)m_1/\sqrt2=n_ea+n_ma_D+Sm_1/\sqrt2=\mathcal{Z}$。]

\begin{proposition}[20.22]
对带一般裸质量 $m_1,\dots,m_F$ 的 $1\le F\le4$ 个基本超多重态，味对称性是 $\mathrm{O}(2F)$ 的 Cartan 子群 $\U(1)^F$。中心荷具有形式
\[
\mathcal{Z}=n_ea+n_ma_D+\sum_{i=1}^{F}S_i\frac{m_i}{\sqrt2},
\]
其中 $S_1,\dots,S_F$ 表示在阿贝尔味对称群 $\U(1)^F$ 下的荷。
\end{proposition}

这表明在绕某个 BPS 粒子变为无质量之点的单值变换下，坐标 $(a,a_D)$ 一般不只经历 $\SL_2(\Z)$ 变

% [原文 p.306]
换，还会获得 $m_i/\sqrt2$ 的整系数线性组合作为可加常数。

\subsection{库仑分支的结构}

$\mathcal{C}^{(F)}$ 的结构可以如 §20.1 那样研究。先考虑 $F=1$，设 $\tilde\tau$ 是 UV 规范耦合，$\tau(u)$ 是 $u\in\mathcal{C}^{(1)}$ 处的低能阿贝尔耦合。

在 $\mathcal{C}^{(1)}$ 的弱耦合区域，即 $|u|\gg|\Lambda_1^2|$ 时，有 $u\simeq a^2$，且
\[
\tau(a)\simeq2\tilde\tau(a)=-\frac{3}{\pi i}\log\Bigl(\frac{a}{\Lambda_1}\Bigr)+\dots,
\]
其中省略号表示瞬子修正。这意味着
\[
a_D=-\frac{3a}{\pi i}\log\Bigl(\frac{a}{\Lambda_1}\Bigr)+\frac{3a}{\pi i}+\dots,
\]
从而给出无穷远处的单值性
\[
(a\ \ a_D)\longmapsto(a\ \ a_D)\begin{pmatrix}-1&3\\0&-1\end{pmatrix}.
\]
现在转向强耦合区域，先设 $m_1\gg\Lambda_1$。在低于 $|m_1|$ 的标度上，有质量超多重态脱耦，有效地给出纯 $4d$ $\mathcal{N}=2$ $\SU(2)$ SYM。后一理论的动力学标度 $\Lambda_0$ 可以在单圈近似下用 $\Lambda_1$ 表出：
\[
\Lambda_0^4\simeq m_1\Lambda_1^3.
\]
% [Tutor Note: 标度匹配核对：$F=1$ 时 $b_1=4-F=3$，纯 SYM（$F=0$）时 $b_1=4$。单圈匹配写作 $\Lambda_{F-1}^{b_1(F-1)}=m_F\Lambda_F^{b_1(F)}$，即 $\Lambda_0^4=m_1\Lambda_1^3$，与正文一致。]
由 §20.1 我们知道 $\mathcal{C}^{(F)}$ 上出现位于 $\pm2\Lambda_0^2$ 的两个奇点：在 $2\Lambda_0^2$ 处，荷为 $(n_e,n_m,S)=(0,1,0)$ 的单极子变为无质量；而在 $-2\Lambda_0^2$ 处，荷为 $(2,1,0)$ 的双子变为无质量。还有第三个奇点在 $u\simeq m_1^2$ 处，BPS 超多重态 $(1,0,1)$ 在那里变为无质量。

现在设 $m_1=0$。量子理论的 $\Z_{12}$ 对称性在 $\mathcal{C}^{(1)}$ 上约化为 $\Z_3$ 对称性。与纯理论一样，$\mathcal{C}^{(1)}$ 不能只在原点有一个奇点，因为 $\operatorname{Im}\tau$ 必须正定。与 $m_1\gg\Lambda_1^2$ 的分析相容的最简单推广，暗示在 $u\propto\zeta_3^k\Lambda_1^2$ 处存在三个奇点，其中 $\zeta_3=\exp(2i\pi/3)$、$k=0,1,2$。由于这些奇点被 $\Z_3$ 置换，相应的单值性彼此共轭。$\mathcal{C}^{(1)}$ 的结构示于图 20.6。

我们直接陈述 $2\le F\le4$ 时库仑分支 $\mathcal{C}^{(F)}$ 的结构，并参见 [314] 与 [339, Chaps. 5, 8, 9] 以了解更多细节：

% [原文 p.307]
\Fig[0.75\linewidth]{figures/fig_20_6.png}{20.6}{带 $F=1$ 个基本超多重态的 $\mathcal{N}=2$ $\SU(2)$ SYM 的库仑分支。}

\begin{itemize}
\item 当 $F=2$ 时，同样的推理给出：裸质量一般时 $\mathcal{C}^{(2)}$ 上有四个奇点。当所有质量为零时，这四个奇点两两合并为两个奇点。这两个奇点被作用在 $\mathcal{C}^{(2)}$ 上的 $\Z_2$ 对称性互换，该对称性来自量子理论的 $\Z_8$ 对称性。
\item 当 $F=3$ 时，对一般的质量取值，库仑分支 $\mathcal{C}^{(3)}$ 含有五个奇点。当所有质量为零时，其中四个奇点合并，剩下两个奇点。此时量子理论的 $\Z_4$ 对称性在库仑分支上平凡地作用。
\item 最后，对 $F=4$，当超多重态裸质量一般时，库仑分支 $\mathcal{C}^{(4)}$ 显示六个奇点。当所有质量为零时，这六个奇点在库仑分支的原点合并。这一情形没有矛盾，因为无穷远处的单值性只表现为电荷共轭，即 $(n_e,n_m)\to-(n_e,n_m)$，等价地 $(a,a_D)\to-(a,a_D)$。此时 $\U(1)_R$ 对称性不反常。
\end{itemize}

\begin{remark}[20.23]
对任何 $1\le F\le4$，取裸质量 $m_F$ 趋于无穷大的极限，得到带 $F-1$ 种味的理论的库仑分支。这一做法提供了确定这些理论库仑分支结构的另一种方法，并可作为所列结果的一个非平凡自洽性检验。
\end{remark}

如 §20.1，可以为 $\mathcal{N}=2$ 矢量超场中的伴随手征超场添加 $\mathcal{N}=1$ 质量项，从而得到 $\mathcal{N}=1$ SQCD 理论。库仑分支除其奇点外被完全提升，而希格斯分支给出 §17.3.5 中描述的模空间。

% [原文 p.308]
\subsection{Seiberg–Witten 曲线}

如 §20.1，由库仑分支分析确定的奇点结构完全决定了全纯（局部）函数 $a(u)$ 与 $a_D(u)$。它们可以定义为复环面 $\Sigma_u$ 上亚纯 $1$-形式 $\lambda$ 的周期，其中 $(\Sigma_u)_{u\in\mathcal{C}^{(F)}}$ 构成 $\mathcal{C}^{(F)}$ 上的全纯复环面族。此外，每个 $\Sigma_u$ 都是与 $u$ 无关的黎曼球面 $\C\PP^1_z$ 的 $2$ 次分歧覆叠。

例如，当 $F=1$ 时，对 $(\Sigma_u,\lambda_u)$ 可由方程
\begin{equation}
(\Sigma_u):\quad \Lambda_1^2z+\frac{2\Lambda_1(x-m_1)}{z}=x^2-u,\qquad
\lambda_u=x\frac{\dd z}{z}
\tag{20.13}
\end{equation}
描述。亚纯 $1$-形式 $\lambda_u$ 要求满足
\[
\frac{\dd\lambda_u}{\dd u}=\omega_u,
\]
其中 $\omega_u$ 是 $\Sigma_u$ 上唯一的（相差归一化）全纯 $1$-形式。该条件连同 Riemann 双线性关系意味着规范耦合的虚部
\[
\tau(u)=\frac{\partial a_D}{\partial a}=\frac{\partial a_D/\partial u}{\partial a/\partial u}
\]
正定，从而 $\dd s^2=\operatorname{Im}\tau(u)\,\dd a\,\dd\bar a$ 是 $\mathcal{C}^{(F)}$ 上良好定义的度量。

\begin{remark}[20.24]
在纯 $\SU(2)$ SYM 中，要求亚纯微分 $\lambda$ 的留数为零。该条件保证特殊坐标 $a$ 与 $a_D$ 的定义只依赖于 $\Sigma_u$ 上圈 $A$ 与 $B$ 的同调类。然而对 $F\ge1$，亚纯微分 $\lambda$ 必须有非零留数，以计入超多重态裸质量对中心荷的贡献。这些留数又允许 $F\ge1$ 时出现的 $a$ 与 $a_D$ 的更一般单值性，如式 (20.12) 所示。例如在式 (20.13) 中，$z=0$ 不是覆叠 $\Sigma_u\to\C\PP^1_z$ 的分歧点，这与纯 SYM 中的情形相反。微分 $\lambda_u$ 在 $z=0$ 于 $\Sigma_u$ 中的原像处有非零留数 $\pm m_1$，这是计入式 (20.12) 的单值性所必需的。
\end{remark}

对定义带 $1\le F\le4$ 种味的理论的 Seiberg--Witten 曲线的方程施加适当的脱耦极限，将重现任何 $F'\le F$ 种味理论的 Seiberg--Witten 曲线。再次参见 [314] 与 [339, Chaps. 5, 8, 9]，以了解带 $1\le F\le4$ 个基本超多重态的四维 $\mathcal{N}=2$ $\SU(2)$ SYM 的 Seiberg--Witten 曲线的计算及其性质的分析。

% [Tutor Note: 用 $m_1\to0$ 核对 (20.13) 与纯 SYM 曲线 (20.7) 的相容性：$m_1=0$ 时 (20.13) 为 $\Lambda_1^2z+2\Lambda_1x/z=x^2-u$；把 $x,z$ 作适当常系数重标度即可化为 (20.7) 的形式，且两族在 $z_\pm$ 相撞时给出的奇点结构同为 $\pm2\Lambda^2$，故脱耦极限自洽。]

% [原文 p.309]
\subsection{带四种基本味的无质量 $\SU(2)$ SYM 的 S 对偶}

对 $F=4$，规范耦合单圈重整化系数为零，意味着耦合在微扰论中是常数。虽然瞬子修正在原则上可以贡献 $\tilde\tau$ 的重整化，但这样的贡献会导致标度不变性与 $\U(1)_R$ 对称性的指数小破缺。这又会使 $\operatorname{Im}\tau\,\dd a\,\dd\bar a$ 失去正定性 [314]。虽然这些问题可能通过额外的非微扰效应得到解决，但简单性与自洽性倾向于假定理论真正标度不变。这是一个广泛接受的观点，以下我们假定如此。

\begin{remark}[20.25]
标度不变性意味着 $\tau=2\tilde\tau_{\mathrm{cl}}$ 与 $u$ 无关，且
\[
a(u)=\sqrt u,\qquad a_D(u)=2\tilde\tau_{\mathrm{cl}}\,a(u)
\]
是精确等式，如同 $\mathcal{N}=4$ SYM 理论中那样。与渐近自由理论相比，这极大地简化了电磁对偶的实现。特别地，BPS 谱是 $\SL_2(\Z)$ 不变的，很像 $\mathcal{N}=4$ SYM 理论。然而与 $\mathcal{N}=4$ SYM 不同，带 $F=4$ 个基本超多重态的 $\SU(2)$ SYM 拥有非平凡的味对称性 $\SO(8)$。这使电磁对偶更丰富；特别地，它与味对称性混合。
\end{remark}

对任何 $F\ge1$，基本超多重态在双子背景中诱导 $2F$ 个费米零模，它们在味对称群 $\SO(2F)$ 的矢量表示中变换。量子化这些零模会把单极子与双子变为 $\SO(2F)$ 的旋量，意味着量子味对称性实际上是 $\SO(2F)$ 的万有覆叠 $\mathrm{Spin}(2F)$。当 $F=4$ 时，$\mathrm{Spin}(8)$ 的中心是 $\Z_2\times\Z_2$，而 $\mathrm{Spin}(8)$ 的表示按 $\Z_2\times\Z_2$ 的元素所标记的类分类。特别地，伴随表示与平凡表示对应于 $o=(0,0)$，矢量表示对应于 $v=(1,0)$，正手征旋量表示对应于 $s=(0,1)$，共轭旋量表示对应于 $c=(1,1)$。重要的是，$\mathrm{Spin}(8)$ 显示\emph{三象性}（triality）：其外自同构群是 $\mathfrak{S}_3$，即三个对象上的置换群，这三个对象在此是矢量表示（定义类 $v$）、旋量表示（$s$）与共轭旋量表示（$c$）。它们都是 $8$ 维的。

\begin{proposition}[20.26]
（[314, Sect. 6]）低能粒子在 $\mathrm{Spin}(8)$ 中心下的荷由它们的电磁荷 $(n_e,n_m)$ 模 $2$ 的约化决定。
\end{proposition}

例如，基本超多重态导致荷为 $(n_e,n_m)=(1,0)$ 的低能 BPS 态，它们在 $\mathrm{Spin}(8)$ 的矢量表示中变换。荷为 $(0,1)$ 的磁单极子反而在 $\mathrm{Spin}(8)$ 的旋量表示中变换。因此，由于 S 对偶交换基本超多重态与磁单极子，它也必须置换 $\mathrm{Spin}(8)$ 的表示。此外，S 对偶意味着对任何给定的互素荷对 $(n_e,n_m)$，存在八个
